\section{Proof-Theoretic Concepts and Basic TFL Metalogic}

So far our proof-theoretic rule system can be summarized as follows. In the starred cases (*), the rule is taken to close a subproof which contains the premises. Similarly, \ephi\eif\epsi is typically shorthand for a subproof proving \epsi from \ephi.

\begin{table}[h]
  \centering
  \label{tab:rules-for-tfl-proofs}
    \begin{tabular}{|c|c|c|}
      \hline
      Rule Group       & \textbf{Introduction} & \textbf{Elimination} \\
      \hline
      \enot            & \enot\ephi,\ephi\eproves\ebot           & * \ephi,\ebot\eproves\enot\ephi  \\
      \eand            &
      \ephi,\epsi\eproves\ephi\eand\epsi &
      \makecell{ \ephi\eand\epsi\eproves\ephi;\\ \ephi\eand\epsi\eproves\epsi }                     \\
      \eor             & \ephi\eproves\epsi &
      \epsi\eor\ephi,\ephi\eif\echi,\epsi\eif\echi\eproves\echi                                     \\
      \eif             & $\ephi\dots\epsi\eproves\ephi\eif\epsi$ &                                  \\
      \eiff            &
      \ephi\eif\epsi,\epsi\eif\ephi\eproves\ephi\eiff\epsi &
      \makecell{\ephi\eiff\epsi,\ephi\eproves\epsi;\\ \ephi\eiff\epsi,\ephi\eproves\ephi }          \\
      \hline
      DeM (\enot\eand) &
      \makecell{ \enot\ephi\eor\enot\epsi \eproves \enot(\ephi\eand\epsi) } &
      \makecell{ \enot(\ephi\eand\epsi) \eproves \enot\ephi\eor\enot\epsi }                         \\
      DeM (\enot\eor)  &
      \makecell{ \enot\ephi\eand\enot\epsi \eproves \enot(\ephi\eor\epsi) } &
      \makecell{ \enot(\ephi\eor\epsi) \eproves \enot\ephi\eand\enot\epsi }                         \\
      DN               & \ephi\eproves\enot\enot\ephi            &  \enot\enot\ephi\eproves\ephi    \\
      DS               & \ephi\eor\epsi,\enot\ephi\eproves\epsi  & n/a \\
      LEM              & \ephi\eif\epsi,\enot\ephi\eif\epsi\eproves\epsi & n/a \\
      \hline
    \end{tabular}
  \caption{Proof Rules for TFL Natural Deduction}
\end{table}

\subsection{Formalizing the Notion of \emph{Proof}}

So far I've been talking informally about valid \emph{arguments}. Something
like $\ephi,\epsi\therefore\ephi\eand\epsi$ outlines the argument, but we
want to be more precise in our discussion of various arguments. First, we
will replace the informal \therefore marker with a more formal meta-logical
operator \eproves and its negation \enotproves.

\begin{definition}[Turnstiles]
  \label{def:turnstiles}
  Let $\ephi_1\dots\ephi_n$ and \epsi be TFL sentences. We can say
  $\ephi_1\dots\ephi_n\eproves\epsi$ iff there is a proof that ends in \epsi
  and has $\ephi_1\dots\ephi_n$ as its premises. We can say
  $\ephi_1\dots\ephi_n\enotproves\epsi$ if there is no such proof.
\end{definition}

\begin{definition}[Proof]
  A TFL \emph{proof} is a sequence of TFL sentences presented in a spefific
  format.

  A proof is indented to represent its subproofs. The unindented portion of the
  proof is called the \emph{main proof} and its column is called the \emph{home
  column}.

  A proof begins with the entry of zero or more \emph{premises}, each entered
  in the home column on its own line along with (i) a line number (ii) the
  premise rule (PR) used as a justification. As a general practice, we will
  write (i) to the left of the TFL sentence and (ii) to its right. A horizontal
  line is drawn underneath the last premise.

  Once the premises are added, we can complete the proof by applying one of the
  following at each step:

  \begin{enumerate}
  \item Remain in the currect column and apply one of the TFL rules to add a new
    sentence on a the next line, giving it the next line number and justifying
    it with the rule name.
  \item Move one column to the right and add a sentence, justifying it with the
    assumption rule (AS) and giving it the next line number. This move opens a
    new \emph{subproof}.
  \item Move one column to the left by applying a rule which \emph{closes} a
    subproof. Give the new line the next line number and justify it using the
    name of the rule applied.
  \item If the most recent line is in the home column, you may end the proof.
    The final line of a proof is its \emph{conclusion}. 
  \end{enumerate}

  where the rules are subject to the following contstraints:

  \begin{outline}[enumerate]
  \1 In order to apply, a rule must take its input:
  \2 from a previous line;
  \2 that does not occur inside a closed subproof.
  \1 On the line immediately after closing a subproof, typically either:
  \2 one of the rules which makes use of that subproof is applied next; or
  \2 a new subproof is started so that we can use a rule which expects multiple
  subproofs. 
  \1 A proof can only end if its final line is in the home column. 
  \end{outline}

  When \eGamma are its premises and \ephi its conclusion, a successful proof
  can be referenced as \eGamma\eproves\ephi. The argument that it proves is
  said to be \emph{valid}. A proof is pure \emph{syntax}. Only the rules of the
  proof system are relevant to its evaluation.
\end{definition}

Of course, as with schemas, one of the things that makes proofs such a powerful
part of logic is the fact that they are \emph{truth-preserving under
substitution of equivalents}. So if \ephi\eproves\echi and \ephi\eiff\epsi,
then \epsi\eproves\echi. Among other things, this means that we can think of a
proof as covering a whole collection of related sentences, namely all those
that can be made from it with consistent substitutions.

\subsection{Proof-Theoretic Concepts}

Now we can use the more precise turnstile notation to express a variety of
\emph{proof-theoretic} concepts. Many of these will be familiar from exploring
analogous \emph{truth-functional} concepts earlier.

First, we note a sentence which can be proven from nothing is a \emph{theorem}.

\begin{definition}[Theorem]
  \label{def:theorem}
  A sentence \ephi is a \emph{theorem} iff \eproves\ephi.
\end{definition}

Similarly, sentences that can be proven from each other are
\emph{interderivable}.

\begin{definition}[Interderivable]
  \label{def:interderivable}
  Two TFL sentences \ephi and \epsi are \emph{interderivable} iff
  \ephi\eproves\epsi and \epsi\eproves\ephi. We can write this as
  \ephi\efollows\eproves\epsi.
\end{definition}

We can generalize the notion of consistency from a comparison of two sentences
to a property of a collection of them.

\begin{definition}[Joint Consistency]
  \label{def:joint-consistency}
  TFL sentences $\ephi_1\dots\ephi_n$ are \emph{jointly inconsistent} iff
  $\ephi_1\dots\ephi_n\eproves\ebot$. Otherwise they are \emph{jointly
  consistent}.
\end{definition}

We can summarize the burdens of proof much as we did for truth tables.

\begin{table}[h]
  \centering
  \begin{tabular}{|c||c|c|}
    \hline
    proving       & \textbf{yes} & \textbf{no} \\
    \hline
    theorem?      & one          & all         \\
    inconsistent? & one          & all         \\
    equivalent?   & two          & all         \\
    \hline
  \end{tabular}
  \caption{Metalogical Burdens of Proof}
\end{table}

In this table we answer the question of how many proofs are needed to show the
property in question. The main takeaway is that, while proving the affirmative
is a straightfoward matter of logic, proving the negative requires metalogical
reasoning about \emph{all possible proofs}.

A special case of joint inconsistency is when there is only one sentence. In
such a case, we can say that the sentence itself is \emph{inconsistent}.

\begin{definition}[Inconsistent Sentence]
  \label{def:tfl-inconsistent-sentence}
  A TFL sentence \ephi is \emph{inconsistent} iff $\ephi\eproves\ebot$.
\end{definition}

Every TFL sentence that isn't inconsistent or a theorem is \emph{contingent}.

\begin{definition}[Contingent Sentence]
  \label{def:tfl-contingent-sentence}
  A TFL sentence \ephi is \emph{(proof-theoretically) contingent} if \ephi is
  neither a theorem nor inconsistent.
\end{definition}

And we can make formal the notion of a \emph{valid argument} discussed
informally earlier as provability of the conclusion from the premises.

\begin{definition}[Valid Argument]
  \label{def:tfl-joint-consistency}
  The argument $\ephi_1\dots\ephi_n\therefore\epsi$ is a \emph{valid} iff
  $\ephi_1\dots\ephi_n\eproves\epsi$.
\end{definition}

\subsection{Proof-theoretic vs. Truth-functional}

From a metalogical perspective, there remains the question of the relationship
between the truth-functional and proof-theoretic approaches to logic.
Metalogical proofs will be a major topic in Logic 2, but we will lay out the
broadstrokes before wrapping up our module on Truth-Functional Logic.

Consider the how the definitions of the major logical concepts vary between
the two systems.

\begin{center}
  \begin{tabular}{c||c|c}
    \bf{concept} & \bf{truth-functional} & \bf{proof-theoretic} \\
    \hline
    Tautology / Theorem &
    always T &
    derived without premises \\
    Contradiction / Inconsistency &
    always F &
    can derive \ebot from it \\
    Contingent &
    sometimes T, sometimes F &
    not a theorem or inconsistency \\
    Equivalent / Interderivable &
    identical main columns &
    derive mutually \\
    Unsatisfiable / Inconsistent &
    no line where all T &
    can derive \ebot \\
    Satisfiable / Consistent &
    some line where all are T &
    cannot derive \ebot from them \\
    Valid argument &
    no lines w F $\ephi_i$ and T \epsi &
    can derive \epsi from $\ephi_1\dots\ephi_n$ \\
  \end{tabular}    
\end{center}

Note that we have introduced four major notational devices that look remarkably
similar. We should pause for a moment to distinguish them all.

\begin{argument}
\item[\eif] The material conditional (\eif) is a \emph{logical} connective. It
  is the only one of these symbols that can be used in a sentence of TFL.
\item[\therefore] The inference marker (\therefore) is an \emph{informal} device
  used in the metalogic to discuss arguments. \eGamma\therefore\ephi does not
  imply that the argument is valid, but does specify its premises (\eGamma)
  and conclusion (\ephi).
\item[\eentails] The entailment marker (\eentails) is a formal device used to
  indicate the \emph{truth-theoretic} validity of an argument.
  $\ephi_1\dots\ephi_n\eentails\epsi$ says that the conclusion (\epsi) is true
  in every situation where the premises ($\ephi_1\dots\ephi_n$) all are.
\item[\eproves] The provability marker (\eproves) is a formal device used to
  indicate the \emph{proof-theoretic} validity of an argument.
  $\ephi_1\dots\ephi_n\eproves\epsi$ says that the conclusion (\epsi) can be
  dervied from the premises ($\ephi_1\dots\ephi_n$) by the natural deduction
  rules for TFL.
\end{argument}

Nevertheless, there are parallels in how the symbols are used. The most
significant parallels (soundness and completeness), taken together, mean that
\ephi\eproves\epsi iff \ephi\eentails\epsi. But first consider a few more
trivial observations.

When we were working on truth tables in Lecture 3, we showed how material
implication and entailment are truth-functionally equivalent (i.e.
\ephi\eentails\epsi iff \eentails\ephi\eif\epsi).

\begin{table}[h]
  \begin{center}
    \begin{tabular}{c c||c||c}
       A& B &  A\eif B &  A\eentails B \\
      \hline
      T    & T    &       T        &          T \\
      T    & F    &       F        &          F \\
      F    & T    &       T        &          T \\
      F    & F    &       T        &          T
    \end{tabular}
  \end{center}
  \label{tab:tt-non-reversability}
  \caption{A truth table showing \ephi\eif\epsi iff \ephi\eentails\epsi}
\end{table}

But note that we have to view the meaning of T and F in the A\eentails B column
in a special way: as \emph{good line} and \emph{bad line}, respectively. Truth
is not an evaluation we make about potential entailments, instead we talk about
validity.

Likewise, we can be explicit about the relationship between \emph{provability}
and \emph{inconsistency}.

\begin{observation}[Proof and Inconsistency]
  Let \eGamma = $\egamma_1\dots\egamma_n$ be a set of TFL sentences and
  \ephi a single TFL sentence. If \eGamma\eproves\ephi then the set that
  results from adding \enot\ephi to \eGamma (\eGamma\ecup\enot\ephi) is
  inconsistent.
\end{observation}

\subsection{Soundness and Completeness}

Two of the metalogical results that we will not prove here are \emph{soundness}
and \emph{consistency}. We will define them and them assume that they are true
as we wrap up our discussion of proofs for logical properties.

Soundness says that we do not get any derivations that are not also entailments.

\begin{definition}[TFL Soundness] % (p 222)
  A proof system is \emph{sound} iff \ephi\eproves\epsi implies
  \ephi\eentails\epsi for all TFL sentences \ephi and \epsi.
\end{definition}

Conversely, completeness says that we do not get any entailments that are not
also derivations.

\begin{definition}[TFL Completeness] % (p 224)
  A proof system is \emph{complete} iff \ephi\eentails\epsi implies
  \ephi\eproves\epsi for all TFL sentences \ephi and \epsi.
\end{definition}

Knowing that a proof system is sound and complete means that we can demonstrate
a logical property using either proof-theoretic or truth-functional appoaches
and know that it holds for both. This means we can leverage the assymmetries in
the last table and approach the problem in whatever way is easiest. The
following table demonstrates such a strategy.

\begin{table}[h]
  \centering
  \label{tab:proof-strategies-for-tfl-metalogic}
  \begin{tabular}{c||c|c}
    \hline    
    \textbf{property} & \textbf{prove present} & \textbf{prove absent} \\
    \hline    
    Tautology, Theorem &
    derive the sentence &
    find a F line \\
    Contradiction &
    derive \ebot &
    find a T line \\
    Contingency &
    find T and F lines &
    prove \ephi or \enot\ephi \\
    Equivalence &
    mutually derive &
    find a line with truth mismatch \\
    Consistency &
    find a universally T line &
    derive \ebot \\
    Validity &
    derive \epsi from $\ephi_1\dots\ephi_n$ &
    find a line with  $\ephi_1\dots\ephi_n$ all T and \epsi F \\
    \hline    
  \end{tabular}
  \caption{Metalogical Proof Strategies for TFL}
\end{table}
